\section{非微扰重正化条件}
\label{sec:idea}

本节详细给出非微扰重正化方案，并讨论为使我们的方法适用所必须满足的条件。

为便于讨论，按紫外行为（当 $a \to 0$ 时）将复合算符分为三大类是方便的：
\begin{enumerate}
\item {\bf 有限算符}：包括矢量流与轴矢流。另一个例子是标量与赝标重正化常数之比 $Z_S/Z_P$。
对这些情形，Ward 恒等式方法同样适用。
\item {\bf 对数发散算符}：这是一大类与强子唯象学相关的算符，包括标量密度与赝标密度、
四费米子 $\Delta F=2$ 算符、能量动量张量的某些分量等。
\item
{\bf 幂次发散算符}：当与低维算符的混合可能时，就会出现这类算符。这发生在具有内禀质量标度的
正规化中，如格点正规化。在格点 $QCD$ 的唯象应用中经常遇到这类例子，例如与 $\Delta I=1/2$
跃迁相关的四费米子算符，或重夸克有效理论（HQET）中的 $1/m$ 阶算符 \cite{effi}--\cite{efff}。
\end{enumerate}

为简单起见，我们先考虑第 2 类的两费米子算符，并在下一节把讨论推广到该类的其他算符以及第 1、3
类算符。因此下面给出的公式可直接应用于赝标密度与标量密度。在整个讨论中，我们假定离散化误差可
以忽略：用场论语言说，这等价于断言重正化格林函数与其在无穷大紫外截断极限下的值没有可观差别。
讨论在小夸克质量极限与欧氏时空中进行。夸克作用的离散化假定采用 Wilson 方式，其特征是拉氏量中
存在显式的手征对称性破缺\footnote{这包括 SW-Clover 作用量 \cite{slami}。}。向交错费米子的推广
是直接的。

考虑两费米子裸格点算符 $O(a)$ 的前向截腿格林函数 $\Gamma_{O}(pa)$，它于动量为
$p$（$p^2=\mu^2$）的离壳夸克态之间、在固定规范（例如朗道规范）中计算。若不固定规范，则在夸克与
胶子外态之间计算的所有格林函数均为零。当规范未完全固定（例如库仑规范）时也是如此。我们通过引入
重正化常数 $Z_O$ 定义重正化算符 $O(\mu)$：
\beq O(\mu)= Z_O\Bigl(\mu a, g(a)\Bigr) O(a). \label{eq:zren} \eeq
$Z_O$ 由施加重正化条件确定：
\beq   Z_O\Bigl(\mu a,  g(a)\Bigr) Z^{-1}_\psi
\Bigl(\mu a,  g(a)\Bigr)
\Gamma_{O}(pa)\vert_{p^2=\mu^2}=1, \label{eq:zdef} \eeq
其中 $Z_\psi$ 是场重正化常数，将在下面定义。这一流程定义的重正化算符 $O(\mu)$ 不依赖于正规化
方案 \cite{altarelli}--\cite{ciu}；但它依赖外部态与规范。这并不影响最终结果：与任意给定微扰论阶
数的重正化条件的连续统计算相结合后，结果将是规范不变且不依赖外部态的。让我们明确式
(\ref{eq:zdef}) 中各个量。$\Gamma_{O}$ 用非截腿格林函数 $G_{O}(pa)$ 与夸克传播子 $S(pa)$ 的期望值
定义\footnote{``期望值''照例指对由蒙特卡罗模拟生成的规范场组态求格林函数的平均，见第
\ref{sec:duef} 节。}：
\beq \Gamma_{O}(pa)= \frac{1}{12} {\rm tr} \left( \Lambda_{O}(pa)\hat  P_O
 \right), \label{eq:proj} \eeq
其中
\beq \Lambda_{O}(pa)=  S(pa)^{-1} G_{O}(pa) S(pa)^{-1}. \label{eq:lam}
\eeq
$\hat P_O$ 是树图级算符上适当的投影算符：对标量（赝标）密度取 $\hat P_O=
\hat 1$（$\hat P_O= \gamma_5$）。因子 $1/12$ 保证迹的正确整体归一化（色$\times$自旋=12）。在定义格林
函数——特别是非微扰情形——时投影算符非常方便，文献 \cite{buras,ciu} 中已广泛使用。当然也可以使用
$Z_O$ 的其他定义。

$Z^{1/2}_\psi$ 是费米子场的重正化常数。它可有不同定义，其中一些在微扰意义上等价。超出微扰论之外，
$Z_\psi$ 最自然的定义来自守恒矢量流 $V^{C}$ 的截腿格林函数。事实上，众所周知对 $V^{C}$ 重正化常数
等于一：
\beq
Z^{-1}_{V^{C}}= \Gamma_{V^C} \times Z^{-1}_{\psi}=
\frac{1}{48} {\rm tr}
\left( \Lambda_{V^{C}_{\mu}} (pa)\gamma_{\mu}\right)\vert_{p^2=\mu^2} \times
Z^{-1}_{\psi}=1,
\eeq
这意味着
\beq
Z_{\psi}=\frac{1}{48} {\rm tr}
\left( \Lambda_{V^{C}_{\mu}} (pa)\gamma_{\mu}\right)\vert_{p^2=\mu^2}.
\label{eq:zpsi}
\eeq

式 (\ref{eq:zdef})--(\ref{eq:zpsi}) 完整地定义了我们的方法。本节余下部分讨论关于其适用性的几个重要
方面。

在上面诸式中，为简单起见我们始终只考虑了前向矩阵元。一般地，人们可以自由地在不同的外动量 $p$ 与
$p^\prime$（$p \neq p^\prime$）处定义重正化条件。夸克态的虚度必须远大于 $\Lambda_{QCD}$。原因在于：
为了得到物理结果，必须把重正化算符 $O(\mu)$ 的矩阵元与一个 Wilson 系数函数相结合；后者在连续统微扰
论中以 $\as$ 在 $\mu$ 量级的标度处展开计算。因此，为使该微扰计算有效，$\mu$ 必须很大。我们方案中的
一个重要问题是：能否在格点上找到一个足够低的标度 $\mu$ 以使 $O(a)$ 效应很小，同时又足够大以使高阶
修正很小。满足上述条件的 $\mu$ 范围依赖于数值模拟的 $g_0^2$ 取值。我们强调：如果这样的窗口在当前格点
模拟中不存在，那么不仅在我们的途径中，而且在任何其他方法中，把格点算符精确匹配到连续统算符都将
不可能。

也许有人预期条件 $\mu \gg \Lambda_{QCD}$ 即足以保证微扰论有效。然而当发生自发对称性破缺时（$QCD$
正是如此），大的 $\mu$ 值可能并不足够，因为存在 Goldstone 玻色子——在我们这里是 π 介子。在低动量转移
$q= p-p^\prime$ 处，格林函数会收到来自 π 极点的无法计算的非微扰贡献。天真的预期是该效应在大 $q^2$ 时
以 $1/q^2$ 消失；然而对费米子场而言，即使 $q^2=0$，该贡献仍正比于 $1/p^2=1/\mu^2$。证明见附录。我们
由此得出结论：大的 $\mu$ 值可以解决此问题。但在低 $\mu^2$ 值处，对赝标密度和轴矢流，我们预期会发现
可见的效应，见第 \ref{sec:casivari} 与第 \ref{sec:nonper} 节。

非微扰方案的另一个危险可能来自 Gribov 副本的存在。此处我们不处理这一问题；文献 \cite{altri,gribov}
已对此作过研究。文献 \cite{gribov} 的结果表明，至少对本研究所关心的量，当前格点模拟中的 ``Gribov
不确定性''至多与统计误差同量级。

可以设想在非常大的 $\beta$ 值下、在物理体积很小的格点上施加重正化条件。这似乎能克服 $\mu$ 窗口存在性
的问题。然而对于物理应用，必须把重正化常数演化到更小的 $\beta$ 值，在大 $\beta$ 下探测不到的非微扰贡献
此时可能变得重要。这说明：对任何方法而言，窗口 $\Lambda_{QCD} \ll \mu \ll 1/a$ 的存在都是必要的，这里
$a$ 是用于物理量数值模拟的格点间距。只有如此，非微扰效应与格点伪迹才能同时保持很小。
